\begin{table*}[htbp]

\centering

\small

% \setlength{\tabcolsep}{1pt}

\renewcommand{\arraystretch}{1.1}

\begin{tabular}{lccccc}

\toprule

\textbf{Setting (SpectraVerse)} & \textbf{R@1} &\textbf{R@5}  & \textbf{R@20} & \textbf{MRR} & \textbf{MCES@1} $\downarrow$ \\

\midrule

\textbf{Pretrained, Frozen, Align} & \textbf{31.327}  & \textbf{41.179} & \textbf{54.061} & \textbf{36.587} & \textbf{7.64} \\

Pretrained, Joint Fine-tune, Contrastive & 27.186  & 36.431 & 50.320 & 32.415 & 7.99 \\

Pretrained, Frozen, Contrastive & 27.186  & 36.431 & 50.320 & 32.415 & 7.99 \\

Random Init, Joint Train, Contrastive & 30.798  & 34.457 & 45.625 & 33.618 & 8.58 \\

Random Init, Joint Train, Align & 24.342  & 31.985 & 41.888 & 28.727 & 10.15 \\

\bottomrule

\end{tabular}

\caption{\textbf{Ablation study on encoder pretraining and training objectives (SpectraVerse test set).}

Results show the impact of pretraining (pretrained vs random initialization), encoder training strategy (frozen vs joint fine-tune), and training objective (align vs contrastive). All models are evaluated on the test set with MCES@1 measuring structural similarity of top-1 predictions.}

\label{tab:ablate_spectraverse}

\end{table*}

